%% kapa: exam=1 round=28 year=2017 period=1 form=A subject=civil_law subject_ko=민법 seq=22 number=22
%% source: Q-Net 감정평가사 기출문제 (cst00309) artlSeq=5199603 file=extracted/1008/5199603/2017년 제28회 감정평가사 1차시험 문제지/1교시 A형.pdf
%% conversion: kordoc 4.24.51 -> tools/gen1.py (auto); review=proofread-sonnet; flags=none
\begin{question}{22}
甲은 乙 소유 토지 위에 식재된 입목등기가 되어 있지 않은 소나무 50그루에 대하여 매매계약 체결과 동시에 소유권을 이전받기로 약정하였다. 甲은 계약 체결 후 잔금을 지급하지 않은 채 乙의 동의 하에 소나무 50그루에 각각 ‘소유자 甲’이라는 표기를 써서 붙였다. 이후 乙은 이 소나무를 丙에게 이중으로 매도하였다. 이에 관한 설명으로 옳은 것은? (다툼이 있으면 판례에 따름)
\begin{choices}
  \item 乙은 여전히 소나무에 대하여 소유권을 가진다.
  \item 甲은 소나무에 대하여 입목등기 없이 소유권을 취득한다.
  \item 丙이 乙과의 계약에 의해 명인방법을 갖추면 丙이 소유권을 취득한다.
  \item 甲은 명인방법을 통해 소나무에 대하여 저당권을 설정할 수 있다.
  \item 甲은 소나무에 대하여 입목등기 없이 丙에게 대항할 수 없다.
\end{choices}
\end{question}
\answer{2}
